\chapter{Strategies for the Rest}
\label{ch:strategies}

A structural analysis owes its readers agency, not just diagnosis. This chapter states what each class of actor should do conditional on this book being broadly right---and notes, where it matters, which moves remain correct even if Chapter~\ref{ch:synthesis}'s falsification branches fire. The organizing principle throughout is the one the precedents teach: \emph{do not defend the rent; position on the commons}.

\section{The United States: win the commoditization race or lose it}

The US strategy since 2019---chokepoint denial plus closed-frontier rents---has produced the outcomes this book documents: each denied input spawned a subsidized substitute, and the rent structure concentrated national exposure into the leveraged trade of Chapter~\ref{ch:trade}. The counterfactual strategy has never been tried: compete \emph{for} the commons rather than against it. Concretely: fund open-weight frontier models as public infrastructure (the gpt-oss reflex, sustained and resourced, rather than one defensive release); make the American open stack---not merely American APIs---the world's default, because the adoption layer, once lost, prices everything above it; and attack the physical constraint that is actually binding---the watt (Chapter~\ref{ch:energy})---with the urgency currently spent on GPU export paperwork. The scarcest honest observation: the US capital structure, not its technology, is the strategic vulnerability---an orderly deflation of the scarcity premium, engineered early, is strictly preferable to defending it to the break.

That instruction assumed commoditizing and denying were separate policies that could be run in parallel. Chapter~\ref{ch:mirror} shows they are not, and this is the sharpest finding in the book for American policy. An export-control regime and an open-weight corporate strategy are aimed at the same object from opposite sides. Open weights released \emph{with their training data} by an American firm lower the cost of closing the model-layer gap for whoever is closing it, and Nemotron ships the datasets and the recipes alongside the weights~\cite{nemotron} [P]. An open instruction set that the same firm has publicly qualified as a CUDA host lowers the cost of the compute-layer detour the controls exist to make expensive~\cite{hc2026-nv-riscv,cuda-riscv} [P]. And a scale-up domain rebuilt out of merchant Ethernet---the majority design at Hot Chips 2026 among parts with an off-package fabric---is precisely the substitution the chokepoint strategy exists to prevent, performed by the firms it exists to protect~\cite{hc2026-scaleup} [P]. None is a concession to China; each is a commercial decision taken against a domestic rival, and each would have been taken in a world with no Chinese AI industry at all. The externality does not care about the motive: a denial policy premised on the denied party having to obtain the controlled capability from the controlling country is undermined as effectively by an American release as by a Chinese one, and the American release arrives with the corpus attached.

\emph{The policy conclusion is not that the firms should stop.} That is the reflex, and it is wrong on this book's own analysis: the same releases are what keep an American stack in contention for the non-aligned buyer of Chapter~\ref{ch:nonaligned}, and a government that suppressed them would be surrendering the adoption layer to preserve a denial instrument whose effectiveness it cannot verify---the solar mistake with a national-security justification attached. The conclusion is narrower and harder. \emph{A denial policy whose effectiveness is contingent on the conduct of firms the government does not control is not a strategy}; it is a hope with a licensing regime attached, and the United States must choose which of its two instruments it means. If it means denial, the controls must reach the conduct that erodes them---which runs immediately into the constitutional and practical difficulty of restraining publication, a difficulty to be priced honestly rather than assumed away, and into the Commerce Department's own finding on the open ISA~\cite{riscv-restrict}. If it means commoditization, the controls should be triaged down to the few places the Chapter~\ref{ch:analogy} model says they buy real time---EUV, HBM process, advanced packaging---and the energy freed spent on the watt and on the open stack. What cannot work is running both at full strength and expecting each to do its job, because in the three instruments above they cancel.

The counter-evidence bounds the finding and deserves the last word. The controls have accomplished something measurable: NVIDIA's Hopper shipments to China ran under one per cent of datacentre revenue in the July 2026 quarter, and the forward outlook assumes \emph{no} China datacentre compute revenue at all~\cite{nvda-segments} [P]. Removing the dominant supplier from the second-largest market is not nothing, whatever else leaks. And the leak is smaller than the argument at full strength implies: the American open models shipped to date take routed volume without taking frontier capability, so what they have lowered is the cost of the volume tier and not the capability tier. The cancellation is real, partial, and incomplete---which is the condition under which a choice can still be made deliberately rather than discovered later in a filing.

\section{Europe: regulation is not production}

Europe enters the fourth war having sat out the previous three's manufacturing and owning the regulatory pen. The trap is repeating solar: tariffs and standards atop a hollowed industrial base. Europe's live assets are ASML (the one true chokepoint, wasting), Mistral and the open-model tradition, EuroHPC's procurement scale, and the credibility to run the certification layer of Chapter~\ref{ch:trust} as a neutral. The strategy that follows: build sovereign inference and smaller sovereign-scale training on whatever open stack is best---which is now a genuine choice rather than a euphemism for Chinese weights (Section~\ref{sec:middle-powers}): accept both, audit both, mirror both; spend the AI Act's institutional energy on the trust stack of Section~\ref{sec:middle-powers}, which the world lacks and Europe is culturally equipped to supply; and convert energy policy into compute policy, because Europe's grid math resembles America's with less gas.

\section{Japan: the third-country playbook, first-person}

Japan's position is unusually concrete because its assets map one-to-one onto this book's elements: frontier-class HPC engineering (the Fugaku-to-FugakuNEXT lineage whose architectural argument Chapter~\ref{ch:compute} carries~\cite{dhm-long}); the world's most defensible remaining semiconductor inputs; memory at the leading edge; and deep bilateral standing in both blocs' technical communities. Chapter~\ref{ch:nonaligned} sets out the record and the liabilities in detail; five moves follow from the analysis. \emph{First}, build the accelerated-CPU line to its conclusion---an FP8-equipped, capacity-rich successor in the LX2/Monaka direction is the machine class this book argues wins the converged AI-plus-simulation workload, and Japan is one of three polities that can execute it [M/S]. \emph{Second}, adopt the open commons---now plural---with the full trust stack of Chapter~\ref{ch:trust}, self-hosted, audited, mirrored and fine-tuned domestically; refusing Chinese weights on principle is strategy-by-flag, and taking an American vendor's weights uninspected because they are American is the same error wearing the other flag. \emph{Third}, bid for the certification-authority vacancy that Section~\ref{sec:middle-powers} prices---convened with partners rather than alone, it is the highest-leverage per-yen investment in the entire space, the FAA role held by a party both blocs can transact with, and Chapter~\ref{ch:consortium} develops its institutional form. \emph{Fourth}, energy realism: Japan cannot win a watts race; it can win the performance-per-watt race, which is where its silicon and systems tradition already lives. \emph{Fifth}, hedge the memory bet both ways---HBM partnership with allied suppliers and LPDDR-capacity architectures domestically, for exactly the dual-track reasons of Chapter~\ref{ch:semi}.

\section{The middle powers: what two commons are worth}
\label{sec:middle-powers}

Europe, Japan, Korea, Canada, Australia, the Gulf and the rest of Chapter~\ref{ch:nonaligned}'s electorate should read Chapter~\ref{ch:mirror} as good news, and the reason deserves stating rather than leaving as an inference. \emph{Two competing commons are a better buyer's market than one}, for the reasons Section~\ref{sec:two-commons} prices row by row. What follows for strategy is not neutrality as a virtue but neutrality as procurement: take from both, run them side by side on one evaluation harness, and let the sponsors compete. The only thing that forfeits the position is exclusive adoption in either direction for reasons of alignment rather than of evidence.

The asset neither commons supplies is the one to build, and it is the asset this book has now named in four chapters: \emph{evaluation and certification}. Both stacks ship free weights and neither ships independent assurance about them, because in both cases the party best placed to certify is the party that paid for the training run---an accelerator vendor defending silicon, or a laboratory operating under an industrial policy. That is a structural vacancy rather than an oversight, and the one layer where a middle power's institutional credibility is the scarce input rather than its capital or its watts. Chapter~\ref{ch:trust} specifies the technical content and Chapter~\ref{ch:consortium} the institutional form; Chapter~\ref{ch:mirror} adds the market size, since the number of artifacts requiring assurance has just doubled while the number of credible assurers has not moved.

And one concrete opportunity has a date on it. The commoditization of the scale-up fabric onto merchant Ethernet materially lowers the cost of a sovereign build: the deepest single source of lock-in three years ago is now switch silicon anyone can buy, the majority design at Hot Chips 2026 among parts with an off-package fabric, with a merchant vendor guiding above \$16B a quarter to supply it~\cite{hc2026-scaleup,broadcom-q2fy26} [P]. A sovereign cluster no longer needs a proprietary interconnect licence to reach defensible scale---which changes the arithmetic of EuroHPC's gigafactory tender and of every national machine specified between now and 2028~\cite{eu-giga-tender}. The date is the next procurement, not a research horizon, and the window shuts when those specifications freeze: a system bought on a proprietary fabric in 2027 carries the dependency for the life of the machine, and the FugakuNEXT tension Chapter~\ref{ch:nonaligned} records is the specific form the mistake takes~\cite{fugakunext}. The counter-argument is real: merchant Ethernet at scale-up is newer, its software less mature, and the firm with the most non-NVIDIA accelerators in the field has just placed its next generation inside NVIDIA's interconnect and memory stack~\cite{aws-nvda} [P]. A buyer who reads that as a reason to wait is not being timid; it is the same evidence read for risk rather than for opportunity, and the falsification test below separates the two readings.

\section{The Global South: the buyers write the ending}

The precedent wars were decided, in the end, by third markets: solar in the desert tenders, EVs in Southeast Asia, telecom in the 4G buildouts. The AI war will be decided by the same buyers, and their rational strategy is uncomfortable for both blocs: take the free stack, keep the exit. Chapter~\ref{ch:nonaligned} states the six moves that follow and they are not restated here; what Chapter~\ref{ch:mirror} adds is that ``the free stack'' is now two stacks with different sponsors and incompatible commercial logics, which strengthens every one of those moves---a buyer with two suppliers of the same free good is a buyer with terms, and the trust-stack insistence that was a condition of adoption is now also a lever in a negotiation. WAICO's 29 founding states have understood this faster than the G7 has~\cite{waico}. Whether the commons they join has one certification authority or two---whether trust becomes a Chinese utility or a neutral one---is the open variable, and it is the same variable Chapters~\ref{ch:trust} and this chapter keep landing on, because it is the war's last genuinely unclaimed high ground.

One correction to the buyer's own picture of the menu belongs here, because the menu is being described to ministries in terms that the record does not support. The buyer is not, at \datafreeze, choosing between a financed Chinese stack and a financed American one. Neither exists in the field: the American programme has seventy-eight proposals and no designated package, the Chinese wrapper has a founding agreement with no compute, financing or datacentre clause, the one Ascend export bid is unsigned and the vendor's output is spoken for at home for a year (Section~\ref{sec:export-offer}). The telecom precedent that the sovereign-stack pitch borrows from ran on sovereign credit and surplus factories, and Chapter~\ref{ch:telecom} records that neither has appeared in AI. So the practical choice a ministry faces this year is narrower and more tractable than the pitch implies: it is between open weights on whatever silicon the buyer can actually obtain---which for most of this chapter's electorate means licensed NVIDIA or AMD parts, or CPU-class hardware, running Chinese or vendor-sponsored Western weights---and a proprietary stack on the same licensed silicon. Framed that way, the decision is the six moves of Chapter~\ref{ch:nonaligned} and not an alignment question at all. And the only executed template is Brazil's: split the programme between vendors so that neither owns the estate, finance it from your own science budget so that neither owns the debt, and write the training of your own people into the contract so that the spillover stays---a structure that, whatever else it does, no single foreign licensing decision can strand and no foreign bank can call~\cite{brazil-ai-supercomputer} [A]. It is also one contract of a quarter-billion dollars, contracted for delivery in July 2027 and not yet delivered, and the Ascend half of it sits under the same US guidance that has been cited to Cairo and Kuala Lumpur; the template is executed, not proven. The buyer that waits for the financed stack to arrive is waiting for an offer that, on both sides, is still a press release.

\section{China: advice for the presumptive winner}
\label{sec:china-advice}

A book titled as this one is owes its subject the same service it gives everyone else: if the analysis is right about the mechanism, it is also right about the mechanism's failure modes, and several are visible \emph{now}, documented largely in Chinese sources. What follows is what Beijing should fix, stated with the same directness as the other sections.

\emph{First, the involution problem is real, recognized at the top, and not yet solved.} The model-layer price war---repeated near-total API cuts cascading through Qwen, Zhipu, MiniMax, and the rest within days of each DeepSeek release---is playbook step~3 (overcapacity as strategy) running \emph{without the step-5 rent-recapture plan attached}. Part~I's precedents survived below-cost phases because the losses purchased learning curves that led somewhere: cells, packs, cars. The model-layer war's own leadership doubts the destination: the State Council has formally warned against ``disorderly competition'' in AI, the pricing-law amendments target involution by name, and Xi Jinping himself has publicly questioned whether every province needs an AI, compute, and EV project~\cite{china-involution}. The fix is not to end competition---the Darwinian dynamic produced DeepSeek, and strangling it is falsification branch (e) of Chapter~\ref{ch:synthesis}---but to move the war's price floor from below-cost tokens to at-cost tokens plus \emph{paid complements}: certification, deployment engineering, vertical solutions. The rent-migration table (Section~\ref{sec:rent-migration}) applies to China too, and Chinese labs are currently giving away row one while building almost nothing in rows two through five.

\emph{Second, the monetization gap is the strategy's soft underbelly.} China's flagship labs generate strikingly little revenue relative to their US counterparts---the post-IPO disclosures of Zhipu and MiniMax show thin revenues and heavy losses, and the domestic enterprise market suffers a documented ``low-trust open-source paradox'' in which Chinese firms themselves hesitate to build on domestic open models~\cite{china-monetization}. Against Anthropic's \$47B run rate, the entire Chinese frontier sector's revenue is a rounding error. Nor is the problem confined to the startups: Alibaba---the family with a billion model downloads and the strongest distribution position in the Chinese ecosystem---reported ex-items net income down 99.7\% and adjusted EBITDA down 84\% in the May 2026 quarter while committing to exceed a three-year, \$56B AI capex plan, with its chief executive stating that ``margin is still secondary''~\cite{alibaba-capex}. That is the American hyperscaler bargain (spend now, monetize later) run by a firm whose model layer earns even less. The state can carry this indefinitely; the question is whether it should. A commons whose builders capture nothing eventually stops attracting the best builders---the talent flywheel runs on equity as well as patriotism, and Hong Kong listings at compressed valuations are a weak substitute for a functioning domestic growth-capital market. The fix is the one Beijing finds hardest: let winners charge. Three moves in the record already point this way and deserve to be read as policy templates rather than defections---Kimi K3's revenue-threshold licence, Alibaba's retention of a metered Max tier above its open families (Section~\ref{sec:maxtier}), and DeepSeek's introduction of peak-hour surge pricing on the cheapest frontier-adjacent model in the world. Rent relocated is the strategy working; rent abolished is the strategy eating itself.

\emph{Third, the compute buildout has repeated the classic overcapacity error, and candor about it would help.} Independent reporting documents hundreds of state-subsidized data centers with utilization far below plan---up to 80\% of some new capacity unused---rental prices collapsing, and the familiar provincial dynamic of building the designated thing regardless of demand~\cite{china-idle-dc}. This book's own busbar casebook (Section~\ref{sec:busbar-casebook}) shows the western hubs working; it does not show all of them being needed. The fix is the one the involution campaign gestures at: consolidate ruthlessly, publish utilization, and let the East-Data-West-Computing framework allocate by measured demand rather than by provincial ambition---because idle capacity is not deterrence, it is depreciation, and the American bear case about Chinese AI (``it's all empty data centers'') gains its plausibility from exactly this.

\emph{Fourth, the trust deficit is the binding constraint on the strategy's global half, and it is self-inflicted where it is not structural.} Chapters~\ref{ch:trust} and~\ref{ch:order66} documented what buyers of the commons must price: censorship layers that transfer into derivatives, provenance disputes, an unowned certification layer. Every quarter the certification vacancy stays unfilled, adoption of Chinese weights in regulated Western and non-aligned sectors stalls at the paying-workload boundary---the C919 pattern. The fix is counterintuitive for a security state but follows directly from the analysis: \emph{sponsor the independent certification of your own artifacts}. Reproducible builds, signed manifests, third-party evaluation access, and behavioral-audit cooperation would cost little, could be led through WAICO or offered to a neutral consortium (Chapter~\ref{ch:consortium}), and would convert the commons' largest liability into a completed sale. A China confident in its artifacts should want them audited; the inference buyers draw from the alternative is priced into every procurement this book has examined.

\emph{Fifth, the physical gaps deserve honesty rather than roadmap theater.} HBM and advanced packaging remain the binding constraints on the accelerator fleet; EUV remains absent; and the analyst consensus puts CXMT roughly three years behind the Korean leaders at the class that matters~\cite{china-hbm-gap}. This book's compute chapters argue China's best move is to \emph{change the game} (capacity-first architectures, the hybrid tier) rather than chase the incumbent's metric---advice this section now returns to its subject: fund the 5\,TB/s knee test and the hybrid tier's placement policy (Appendix~\ref{app:lpddr}) at national priority, because they are cheap, decisive, and strategy-defining, and stop announcing HBM schedules that the record keeps missing, because each slipped roadmap feeds the skeptics' case at home and abroad.

\emph{Sixth, submit to the scoreboards---and finish the auditability you have started.} No Chinese-designed accelerator has ever appeared in MLPerf Training; no LX2-class machine has published the AI proof suite; the embodied fleet publishes shipments, not intervention-free hours. The model layer shows both halves of the problem at once, and the good half deserves crediting first: Qwen3.8-Max shipped with a release blog carrying the full benchmark table \emph{and a complete project trace on GitHub} for its flagship autonomy demonstration~\cite{qwen-announcement}---which is more verification apparatus than most Western frontier claims carry, and exactly the practice this book has been asking for. The gap is what sits beside it: no architecture report, no activated-parameter disclosure, no training-hardware statement, no hallucination rates of the kind the prior generation published, and---most consequentially---an autonomous silicon-design claim (Section~\ref{sec:offshore}) that is potentially the most strategically important assertion any Chinese lab has made, shipped without any artifact a third party could check. The recommendation follows directly: publish the design flow's audit trail the way the coding run's was published. A verified AI-designed part would do more for China's position than any benchmark table, because it is the claim the West would most like to disbelieve and the one that, if true, dissolves the chokepoint strategy entirely. Every unaudited claim is discounted by the buyers who matter; this particular one is being discounted at a rate that costs China real strategic credit. The pattern---dominate the metrics you control, avoid the ones you do not---is rational for a challenger and corrosive for an incumbent, and China is becoming an incumbent. The strategy that wins the next decade is the one this book has attributed to China at its best: \emph{show the receipts and let the cost curve argue}. It worked for solar. It will work here, if the receipts are real---and if they are not, better to learn it before the world does.


\emph{Seventh, and newest: the principal move is now being run against you, and the advice above assumes it is not.} Every recommendation above presumes that commoditizing the model layer is China's instrument and the West's problem. Chapter~\ref{ch:mirror} records the year that stopped being true: an American accelerator vendor now releases open weights with their training data, has bought a training platform and 109 of the researchers who ran it, and lists the open-model ecosystem among its own demand drivers~\cite{nemotron,nvda-poolside,nvda-huang-q2} [P]. If the West commoditizes that layer competently, China's advantage there compresses---not because Chinese models get worse, but because free stops being a differentiator when the artifact beside it is free and ships its corpus too. The consequence is a reallocation rather than a defeat: China's durable edge moves to the layers the West is conspicuously \emph{not} commoditizing---firm deliverable power at the busbar (Chapter~\ref{ch:energy}), the deployment-generated data only a country running the machines at population scale accumulates (Chapter~\ref{ch:data}), physical production and the embodied fleet (Chapter~\ref{ch:embodied}), and the domestic silicon floor no foreign decision can withdraw (Chapter~\ref{ch:semi}). Those are the rows to fund, and three of the recommendations above---consolidate the idle capacity, let the winners charge, submit the artifacts---are what makes funding them affordable. The counter-evidence belongs in the same paragraph: a West that commoditizes the model layer also \emph{validates} the open-weight strategy and enlarges the ecosystem in which China is the larger supplier, and in the last week of August three of the four significant open-weight releases were still Chinese~\cite{hunyuan-hy4} [V]. The compression is prospective rather than observed---which is exactly the kind of thing a presumptive winner notices too late.

\section{The common instruction}

Reduced to one line per actor: the US should choose between its two instruments and, having chosen, commoditize before being commoditized; Europe should certify what it cannot manufacture; Japan should build the machine class the war is converging on and referee the trust layer; the middle powers should take from both commons and certify what neither will; the Global South should take the commons and unionize the demand side; and China should let its winners charge, consolidate its overcapacity, and submit its artifacts to scoreboards it does not control. For every actor except the two principals, those instructions converge on one institution, which Chapter~\ref{ch:consortium} specifies. And all of them share one instruction the precedents make non-negotiable: read the indicator table of Chapter~\ref{ch:trade} quarterly, because every strategy above has a branch point there, and the actors who moved early in solar, batteries, and EVs are the only ones the record remembers kindly.

Which of these strategies is right is not a matter of taste, and the register in Appendix~\ref{app:dashboard} carries the two tests that settle it. If merchant-GPU datacentre revenue growth decelerates while custom-ASIC revenue holds the trajectory Broadcom has guided, the commoditization of the silicon layer is working, the pincer of Chapter~\ref{ch:mirror} is real, and the advice above is calibrated correctly: the American choice becomes urgent, the middle powers' sovereign-build window is open, and China's advantage compresses at the model layer while holding at the physical ones~\cite{broadcom-q2fy26}. If instead custom silicon keeps converging onto NVIDIA's interconnect and memory---as Trainium4 has already done, adopting NVLink Fusion and NVIDIA custom HBM~\cite{aws-nvda} [P]---then the wave is being absorbed, the accelerator vendor priced its giveaway of the model layer correctly, and a sovereign build on merchant parts is a more expensive bet than it looks. Both tests resolve by the end of 2027---inside the planning horizon of every actor in this chapter and outside the attention span of most of them.
